\section{USAJMO 2025}
        \begin{enumerate}
        	\item (\textit{USAJMO 2025}) Let $\mathbb Z$ be the set of integers, and let $f\colon \mathbb Z \to \mathbb Z$ be a function. Prove that there are infinitely many integers $c$ such that the function $g\colon \mathbb Z \to \mathbb Z$ defined by $g(x) = f(x) + cx$ is not bijective.
Note: A function $g\colon \mathbb Z \to \mathbb Z$ is bijective if for every integer $b$, there exists exactly one integer $a$ such that $g(a) = b$.


 



\textbf{Solution}

Note that $c=f(a)-f(a+1)$ makes $g(a)=g(a+1),$ so for all $c=f(a)-f(a+1)$ we have that $g(x)$ is not bijective. If there are infinitely many possible values of $c,$ we are done. Otherwise, there are only finitely many possible values of $f(a)-f(a+1)$ so there are only finitely many possible values of $-(f(a)-f(a+1))=f(a+1)-f(a).$ This means there is a minimum $d$ of $f(a+1)-f(a).$ It is clear for $x>0$ that
 \[f(x)=f(0)+\sum_{a=0}^{x-1}(f(a+1)-f(a))\ge dx+f(0)\]and for $x<0$ that
 \[f(x)=f(0)-\sum_{a=x}^{-1}\le dx+f(0).\]Then, let $c>-d+3$ be an integer. We have for integers $x>0,$ $x\ge 1,$ so
 \[g(x)=f(x)+cx\ge (d+c)x+f(0)>3x+f(0)>3+f(0),\]if $x=0$ we have $g(x)=f(0),$ and if $x<0$ we have
 \[g(x)=f(x)+cx\le (d+c)x+f(0)<3x+f(0)<f(0).\]This means that there is no $a$ such that $g(a)=f(0)+1,$ so $g(x)$ is not bijective.$\blacksquare$

 \textbf{Proof 2} Take $c = f(a) - f(a+1)$ and $g(a) = g(a+1)$, which is not injective, which implies $g$ is not bijective. We provide another approach.

 Let there exist infinitely many $c$. This is possible if one considers all pairs $(a,a+1)$, so we are done.

 If there are only finitely many $c$, then we prove $g$ cannot be surjective, and therefore not bijective. That is, all sets of $c \in \{d_{\min}, \dots, d_{\max} \}$. Choose a $c = M+1$, where $M = d_{\max}$. Then,
 \[g(a+1) - g(a) = f(a+1)-f(a) + c\]
 \[=-d_{\max} + c\]
 \[\geq 1\]

 which is viable. However, if $c > M+1$,

 
 \[g(a+1) - g(a) = f(a+1)-f(a) + c\]
 \[=-d_{\max} + c\]
 \[\geq 2\]

 Which is indeed not surjective and thus $g$ is not bijective.

 The proof is complete. $\square$

 


	\item (\textit{USAJMO 2025}) Let $k$ and $d$ be positive integers. Prove that there exists a positive integer $N$ such that for every odd integer $n>N$, the digits in the base-$2n$ representation of $n^k$ are all greater than $d$.


 



\textbf{Solution 1}

We define a remainder operation $\,a \bmod b\,$ to be the remainder when $a$ is divided by $b$. Also, let $\lfloor x\rfloor$ be the usual floor function.

 Base-$(2n)$ Representation:
 \[n^k \;=\; a_{k-1}\,(2n)^{k-1} \;+\; a_{k-2}\,(2n)^{k-2} \;+\;\dots\;+\;a_1\,(2n)\;+\;a_0,\]where each $a_i$ satisfies $0 \le a_i < 2n.$  Hence, the base-$(2n)$ representation of $n^k$ is $a_{k-1}\,a_{k-2}\,\dots\,a_1\,a_0.$

 Leading Digit:
 \[a_{k-1}  \;=\;  \left\lfloor  \dfrac{n^k}{(2n)^{k-1}} \right\rfloor \;=\; \left\lfloor \dfrac{n}{2^{k-1}} \right\rfloor.\]

 General Digit Formula:For $0 \le i < k,$
 \[a_i  \;=\; \left\lfloor  \dfrac{\,n^k \bmod (2n)^{\,i+1}}{(2n)^i} \right\rfloor.\]

 Because $n$ is odd, one can show 
 \[n^k \bmod (2n)^{\,i+1} \;\ge\; n^{\,i+1},\]which implies
 \[a_i  \;\ge\; \left\lfloor \dfrac{n^{\,i+1}}{(2n)^i} \right\rfloor \;=\; \left\lfloor \dfrac{n}{2^i} \right\rfloor \;\ge\; 2^{\,k-1-i} \left\lfloor \dfrac{n}{2^{\,k-1}} \right\rfloor.\]

 As $n$ grows large, $\bigl\lfloor n / 2^{\,k-1}\bigr\rfloor$  becomes arbitrarily big, so each digit $a_i$ eventually exceeds any fixed $d.$  Hence there exists an integer $N$ such that for all odd $n > N,$ the digits $a_i$ in the base-$(2n)$ representation of $n^k$ are all $> d.$

 


	\item (\textit{USAJMO 2025}) Let $m$ and $n$ be positive integers, and let $\mathcal R$ be a $2m\times 2n$ grid of unit squares.


 A domino is a $1\times2$ or $2\times1$ rectangle. A subset $S$ of grid squares in $\mathcal R$ is domino-tileable if dominoes can be placed to cover every square of $S$ exactly once with no domino extending outside of $S$. Note: The empty set is domino tileable.


 An up-right path is a path from the lower-left corner of $\mathcal R$ to the upper-right corner of $\mathcal R$ formed by exactly $2m+2n$ edges of the grid squares.


 Determine, with proof, in terms of $m$ and $n$, the number of up-right paths that divide $\mathcal R$ into two domino-tileable subsets.


 



\textbf{Solution 1}

We claim the answer is $\boxed{\binom{m+n}{m}^2}$. Color $\mathcal R$ in a black-and-white checkerboard pattern such that the lower left square is black. Suppose an up-right path $P$ splits $\mathcal R$ into two subsets $S$ and $S'$ such that $S$ is below $P$. Call a subset balanced if there are an equal number of black and white squares.

 Claim 1: $S$ and $S'$ are tileable iff $S$ is balanced.

 Proof: Since $\mathcal R$ itself is balanced and rotating it by $180^\circ$ swaps $S$ and $S'$, it suffices to only consider $S$. Note that dominoes each cover exactly one black and one white square, so $S$ being tileable implies $S$ is balanced.

 We will proceed with the converse by attempting to remove a domino from $S$ such that $S$ can still be traced by a new up-right path. Indeed, so long as $P$ contains a sequence of $\uparrow\rightarrow\rightarrow$ or $\uparrow\uparrow\rightarrow$, we can replace these with $\rightarrow\rightarrow\uparrow$ and $\rightarrow\uparrow\uparrow$, respectively, and effectively remove a domino (this also preserves balanced-ness). Repeating this until we reach the empty set suffices, as we can reconstruct using the sequence of domino removal. The only cases where this fails is when $P$ forms a staircase and remains on the upper or left edges of $\mathcal R$ for at most one edge each. However, we can trivially observe that $S$ is then unbalanced, since if the largest diagonal in $S$ has $k$ squares (which must be of the same color), $k+(k-2)+\cdots>(k-1)+(k-3)+\cdots$. $\square$

 Assign coordinates to the gridline intersections such that the lower left corner is $(0,0)$ and the upper right is $(2m,2n)$. Assign each $\rightarrow$ move with starting point $(x,y)$ a two-letter designation, where the first letter is $o$ or $e$ depending on the parity of $x+y$ and the second letter depends on the parity of $x$. (I blame this naming system on leo; he was making some really weird sounds :P)

 Claim 2: $S$ is balanced iff there are an equal number of $\rightarrow$'s with starting letter $o$ and $e$.

 Proof: Consider $S$ as a union of multiple columns of squares, each topped with a $\rightarrow$. Upon inspection, $oe$'s give one more black square than white square, while $eo$'s give one more white square. There are an even number of squares, and therefore equal number of black and white squares, under an $oo$ or $ee$. Thus, for $S$ to be balanced, there must be an equal number of $oe$'s and $eo$'s; but since there are also an equal number of odd-numbered and even-numbered columns, we must have $oe+ee=eo+oo$, so $ee=oo$ and $oe+oo=eo+ee$, as desired. Reversing the logic gives the converse. $\square$

 We finish by splitting moves into two sets based on the parity of $x+y$ and then ordering $m~\rightarrow$'s amongst the $m+n$ moves in each set, giving the desired $\tbinom{m+n}{m}^2$ paths. $\square$

 



\textbf{Solution 2}

$\boxed{\binom{m+n}{m}^2}$

 $\textbf{Proof.}$View the grid as lattice points $(x,y)$ with $0\le x\le 2n$ and $0\le y\le 2m$.Partition the $2m\times 2n$ rectangle into $mn$ disjoint $2\times 2$ blocks
 \[B_{i,j}=\{2i,2i+1\}\times\{2j,2j+1\}, \qquad 0\le i\le n-1,\ 0\le j\le m-1,\]where $(x,y)$ denotes the unit square with lower-left corner $(x,y)$.

 Let $P$ be an up-right path from $(0,0)$ to $(2n,2m)$.Inside each $2\times 2$ block $B_{i,j}$, the restriction of $P$ is a monotone lattice pathfrom $(2i,2j)$ to $(2i+2,2j+2)$. There are six such possibilities; exactly two of themturn at the block center $(2i+1,2j+1)$.

 $\textbf{Claim 1.}$If $P$ turns at $(2i+1,2j+1)$ for some $i,j$, then one of the two regions cut by $P$is not domino-tileable.

 $\textit{Proof.}$A turn at the center of a $2\times 2$ block forces exactly $1$ or $3$ of the $4$ unit squaresof that block to lie on one side of $P$. Any domino tiling covers an even number of squares,so that side cannot be tiled. $\square$

 $\textbf{Claim 2.}$If $P$ never turns at any $(2i+1,2j+1)$, then both regions cut by $P$ are domino-tileable.

 $\textit{Proof.}$If no block-center turn occurs, then in each $2\times 2$ block the path leaves $0$, $2$, or $4$squares on each side, and when $2$ squares occur they share an edge.Hence each side can be tiled independently within each block(by $0$, $1$, or $2$ dominoes), giving a tiling of each entire region. $\square$

 By Claims 1 and 2, the desired paths are exactly those with no turns at block centers.

 Now group the rows into $m$ strips of height $2$: strip $j$ consists of rows $y=2j$ and $y=2j+1$.For such a path $P$, let $r_{j,0}$ be the number of right-steps taken along $y=2j$,and let $r_{j,1}$ be the number of right-steps taken along $y=2j+1$.Avoiding block-center turns forces both $r_{j,0}$ and $r_{j,1}$ to be even. Write
 \[r_{j,0}=2a_j,\qquad r_{j,1}=2b_j\]with $a_j,b_j\ge 0$. Since $P$ has exactly $2n$ right-steps,
 \[\sum_{j=0}^{m-1}(2a_j+2b_j)=2n,\]so
 \[\sum_{j=0}^{m-1}a_j+\sum_{j=0}^{m-1}b_j=n.\]

 The sequence $(a_0,\dots,a_{m-1})$ uniquely encodes an up-right path from $(0,0)$ to $(n,m)$by taking $a_0$ right-steps, then an up-step, then $a_1$ right-steps, and so on.Similarly, $(b_0,\dots,b_{m-1})$ encodes another up-right path from $(0,0)$ to $(n,m)$.

 Conversely, given any ordered pair of up-right paths in an $m\times n$ grid,doubling every right-step and placing one path on even rows and the other on odd rowsproduces a unique up-right path in the $2m\times 2n$ grid with no block-center turns.Thus there is a bijection between the desired paths and ordered pairs of up-right pathsin an $m\times n$ grid.

 The number of up-right paths from $(0,0)$ to $(n,m)$ is $\binom{m+n}{m}$.Therefore the number of desired paths is $\binom{m+n}{m}^2$.

 


	\item (\textit{USAJMO 2025}) Let $n$ be a positive integer, and let $a_0,\,a_1,\dots,\,a_n$ be nonnegative integers such that $a_0\ge a_1\ge \dots\ge a_n.$ Prove that

 \[\sum_{i=0}^n i\binom{a_i}{2}\le\frac{1}{2}\binom{a_0+a_1+\dots+a_n}{2}.\]
Note: $\binom{k}{2}=\frac{k(k-1)}{2}$ for all nonnegative integers $k$.


 



\textbf{Solution 1}

By Vandermonde's,
 \[\frac 12\binom{\sum a_i}2=\frac 12\left(\sum\binom{a_i}2+\sum_\text{sym}a_ia_j\right)\ge\frac 12\sum ia_i^2\ge\sum i\binom{a_i}2,\]with equality at $a_0=0,1$ and $a_i=0$ for $i>0.~\square$

 



\textbf{Solution 2}

We proceed by induction.

 Base case: $n=1$. Note that $\frac{1}{2}\binom{a_0+a_1}{2} \ge \frac{1}{2}\binom{a_1+a_1}{2} = 2a_1(2a_1-1)/4 \ge a_1(a_1-1)/2 = \binom{a_1}2$

 Inductive Hypothesis: Assume $\sum_{i=0}^n i\binom{a_i}{2}\le\frac{1}{2}\binom{a_0+a_1+\dots+a_n}{2}$

 Inductive Step: We try to show it works for $n+1$: $\sum_{i=0}^{n+1} i\binom{a_i}{2}\le\frac{1}{2}\binom{a_0+a_1+\dots+a_{n+1}}{2}$

 By the Inductive Hypothesis all we need to show is $(n+1)\binom{a_{n+1}}{2}\le \frac{1}{2}\binom{a_0+a_1+\dots+a_{n+1}}{2}-\frac{1}{2}\binom{a_0+a_1+\dots+a_{n}}{2}$

 Let $s=a_0+a_1+..+a_n$ and $c=a_{n+1}$. Now we need $(n+1)\binom{c}{2}\le \frac{1}{2}\binom{s+c}{2}-\frac{1}{2}\binom{s}{2}$. The LHS is $\frac{1}{2}(n+1)(c)(c-1)$ the RHS is $\frac{1}{4}(2sc+c^2-c)$. If $c=0$ we're done, otherwise we divide by $c$ and multiply by $4$ to get $2(n+1)(c-1)=2nc+2c-2n-2$ which we are trying to show is $\le 2s+c-1$. Note that $s\ge (n+1)c$ by the problem statement condition. So it simplifies to $2nc+2c-2n-2 \le 2nc+3c-1$ which is obvious.

 


	\item (\textit{USAJMO 2025}) Let $H$ be the orthocenter of acute triangle $ABC$, let $F$ be the foot of the altitude from $C$ to $AB$, and let $P$ be the reflection of $H$ across $BC$. Suppose that the circumcircle of triangle $AFP$ intersects line $BC$ at two distinct points $X$ and $Y$. Prove that $C$ is the midpoint of $XY$.


 



\textbf{Solution 1}

Let AP intersects BC at D. Extend FC to the point E on the circumcircle $\omega$ of $FAP$. Since $H$ is the orthocenter of $\Delta ABC$, we know that $HD = DP$ or $HP = 2HD$, and $AH \cdot HD = CH \cdot HF$. Next we use the power of H in $\omega$: $AH \cdot HP = CH \cdot HE$. These relations imply that $HE = 2HF$.

 Hence $C, D$ are midpoints of $HE, HP$ respectively. By midline theorem, $CD // EP$. Since $AD \perp CD$, we have $AD \perp EP$. This implies that $\angle APE = 90^{\circ}$. Consequently, $AE$ is the diameter of $\omega$. Let $G$ be the midpoint of $AE$ which is also the center of $\omega$. $G,C$ are midpoints of $AE, EH$ respectively. By the midline theorem again, we have $GC//AH$, consequently, $GC \perp BC$. This implies that $GC$ is the perpendicular bisector of the chord $XY$ hence $C$ is the midpoint of $XY$.    ~ Dr. Shi davincimath.com

 



\textbf{Solution 2}

Denote $O_1$ as the center of $(ABC)$, $O_2$ as the center of $FAP$, $K$ as the midpoint of $AF$, $M$ as the midpoint of $AC$, and $N$ as the midpoint of $AP$. It suffices to show that $\angle{O_2CB}=90$.

 Claim: $O_1MO_2C$ is cyclic. 

 Proof: Since $AK=FK$ and $AM=MC$, KM is a midline of $\triangle{AFC}$ and $KM\parallel FC$. $KO_2\parallel FC$ as well since $\angle{AKO_2}=\angle{AFC}=90$, so $M$ lies on $KO_2$. Next, note that $P$ lies on $(ABC)$, so the perpendicular bisector of $AP$ through $N$ passes through $O_1$. In other words, $N, O_1$, and $O_2$ are collinear. Since $NO_2$ and $BC$ are both perpendicular to $AP$, it follows that they are parallel. Since $KO_2\parallel FC$ and $NO_2\parallel BC$, then $\angle KO_2N=\angle{FCB}$.Finally, we have that 
 \[\angle{MO_2O_1}=\angle{KO_2N}=\angle{FCB}=90-B=\angle{MCO_1},\]and thus $O_1MO_2C$ is cyclic. It follows that $\angle O_1O_2C=\angle{O_1MC}=90$, so $\angle{O_2CB}=180-\angle{O_1MC}=90$, as desired.

 -mop

 



\textbf{Solution 3}

Connect $HP$ and have $HP$ intersect $XY$ at $W$. Also extend $FC$ past point $C$ and have it intersect with the circle at point $D$.

 Since $P$ is the reflection of $H$ over $BC$, we know that $HP\perp XY$. Since $H$ is the orthocenter, we can draw the altitude and tell that $A$, $H$, and $P$ are collinear.We know $m\angle{AFH} = m\angle{CWH}=90^{\circ}$ and $m\angle{FHA} = m\angle{WHC}$, so $\triangle{AHF} \sim \triangle{CHW}$ by AA, so $m\angle{FAH} = m\angle{WCH}$.

 $m\angle{FAP} = \frac{1}{2}m\overarc{FP} = \frac{1}{2}(m\overarc{FX} + m\overarc{XP})$ and $m\angle{FCX} = \frac{1}{2}(m\overarc{FX} + m\overarc{DY})$. From this, we can tell that $m\overarc{XP} = m\overarc{DY}$. Therefore, $m\overarc{XD} = m\overarc{PY}$ and $XD = PY$.

 If we connect $HY$, we can tell that that $PY = HY$ due to $P$ being the reflection of $H$ and $WY$ being perpendicular to $HP$, so $XD = HY$. In addition, $m\angle{HYW} = m\angle{PYW} = \frac{1}{2} m\overarc{XP} = \frac{1}{2} m\overarc{DY} = m\angle{YXD}$. Also, $m\angle{HCY} = m\angle{XCD}$ because they are vertical angles.

 So, $\triangle{HCY} \cong \triangle{XCD}$ because of SAA. From this we can conclude that $XC = CY$, so $C$ is the midpoint of $XY$.

 


 



\textbf{Solution 4}

Let $D$ be the foot of the altitude from $A$ to $BC.$ By Power of a Point, we have 
 \[BF\cdot BA = BX \cdot BY = (BC-CX)(BC+CY) = BC^2 + BC(CY-CX) - CX\cdot CY\]and 
 \[DB\cdot DC = DX\cdot DY = (CX-CD)(CD+CY) = CX\cdot CY - CD(CY-CX) - CD^2.\]Adding, we get
 \[BF \cdot BA + DB\cdot DC = BC^2-DC^2 + (BC-CD)(CY-CX).\]It is well known that the reflection of $H$ over $AB,$ which we denote by $P,$ lies on $(ABC).$ Then, let $\angle{BAP} = \angle{BCP} = \angle{BCH} = \theta.$ We have 
 \[BF\cdot BA + DB\cdot DC = (BC\sin\theta)\cdot BA + DB\cdot DC = BC\cdot (BA \sin\theta) + DB\cdot DC = DB\cdot (BC+DC) = (BC-DC)(BC+DC) = BC^2-DC^2.\]Thus, $(BC-CD)(CY-CX)=0,$ and since $BC\neq CD$ we have $CY=CX.$ Hence, $C$ is the midpoint of $XY.$~TThB0501

 



\textbf{Solution 5}

Let Q be the antipode of B. Claim — AHQC is a parallelogram, and AP CQ is an isosceles trapezoid. 

 Proof. As AH $\perp$ BC $\perp$ CQ and CF $\perp$ AB $\perp$ AQ. Let M be the midpoint of QC. 

 Claim — Point M is the circumcenter of triangle AFP. 

 Proof. It’s clear that MA = MP from the isosceles trapezoid. 

 As for MA = MF, let N denote the midpoint of AF; then MN is a midline of the parallelogram, so MN $\perp$ AF. Since CM $\perp$ BC and M is the center of (AF P), it follows CX = CY .

 


	\item (\textit{USAJMO 2025}) Let $S$ be a set of integers with the following properties:


 $\bullet$ $\{ 1, 2, \dots, 2025 \} \subseteq S$.


 $\bullet$ If $a, b \in S$ and $\gcd(a, b) = 1$, then $ab \in S$.


 $\bullet$ If for some $s \in S$, $s + 1$ is composite, then all positive divisors of $s + 1$ are in $S$.


 Prove that $S$ contains all positive integers.


 



\textbf{Solution 1}

We will show that if $\{1, 2, 3, \dots, n\}\in S$, then $n+1\in S$. Note that if $n+1$ is composite, we are done, so we assume that $n+1$ is an odd prime.

 Case 1: $n+1\ne 2^k\pm 1$ for some integer $k>2$.Since $n+2\ne 2^k$, we have that $2^{v_2 (n+2)}, \frac{n+2}{2^{v_2 (n+2)}}\le n\implies 2^{v_2 (n+2)}, \frac{n+2}{2^{v_2 (n+2)}}\in S\implies n+2\in S$. Either $v_2(n)=1$ or $v_2(n+2)=1$, so $2^{v_2(n)+v_2(n+2)}=2^{1+\max(v_2(n), v_2(n+2))}=2\cdot 2^{\max(v_2(n), v_2(n+2))}\le 2\cdot \frac{n+2}{3}\le n\implies 2^{v_2(n)+v_2(n+2)}\in S$. Since $\gcd(n, n+2)=2$, we have that $\gcd\left(\frac{n}{2^{v_2(n)}}, \frac{n+2}{2^{v_2(n+2)}}\right)=1$, so 
 \[2^{v_2(n)+v_2(n+2)}\cdot \frac{n}{2^{v_2(n)}}\cdot \frac{n+2}{2^{v_2(n+2)}}=n(n+2)\in S\implies n(n+2)+1=(n+1)^2\in S\implies n+1\in S.\]

 Case 2: $n+1=2^k+1$ for some integer $k>2$.Note that $k$ must be even, otherwise $2^k+1\equiv 0\pmod 3$. Then we have the following:
 \[\left(2^{\frac{k}{2}}+1\right)^2-1=2^{\frac{k}{2}}\left(2^{\frac{k}{2}}+2\right)=2^{\frac{k}{2}+1}\left(2^{\frac{k}{2}-1}+1\right)\in S \implies\left(2^{\frac{k}{2}}+1\right)^2\in S\]
 \[\implies 2^{\frac{k}{2}-1}\left(2^{\frac{k}{2}}+1\right)^2 = \left(2^{\frac{k}{2}-1}+1\right)\left(2^k+1\right)-1\in S\implies \left(2^{\frac{k}{2}-1}+1\right)\left(2^k+1\right)\in S\implies 2^k+1\in S,\]as desired.

 Case 3: $n+1=2^k-1$ for some integer $k>2$.Let $x<k$ be a positive integer. Then 
 \[(2^x-1)(2^k-1)-1=2^x\left(2^k-2^{k-x}-1\right)\in S\implies 2^k-1\in S,\] as desired.

 Hence, since $\{1, 2, 3, \dots, n\}\in S\implies n+1\in S$, the set $S$ spans all positive integers.

 -mop

 


\end{enumerate}